\begin{abstract}
We prove that every irreducible integer quartic with no fixed prime-square
divisor takes squarefree values with the predicted positive Euler-product
density. More generally, we obtain the corresponding $(d-2)$-power-free
density in degrees $4\le d\le8$. The proof combines number-field factorization,
determinant estimates with adaptive auxiliary primes, and explicit
low-degree geometry. Together with Browning's theorem for higher degrees,
this gives the $(d-2)$-power-free density for every $d\ge4$.
\end{abstract}

\section{Introduction}
\label{sec:introduction}

For an integer $k\ge2$, an integer $a$ is \emph{$k$-power-free} if no prime
power $p^k$ divides $a$. This convention includes negative integers and
excludes zero. Given $f\in\Z[x]$, write
\[
 \rho_f(q)=\#\{a\in\Z/q\Z:f(a)\equiv0\pmod q\},\qquad
 S_{f,k}(X)=\#\{1\le n\le X:f(n)\text{ is $k$-power-free}\}.
\]
The necessary local condition is
\begin{equation}
 \rho_f(p^k)<p^k\quad\text{for every prime }p.
 \label{eq:local-admissibility}
\end{equation}
It says that no prime $k$th power divides every value of $f$. The natural
candidate for the density is the product of the proportions that survive
these separate congruence conditions.

\begin{theorem}\label{thm:density}
Let $f\in\Z[x]$ be irreducible over $\Q$, of degree $d$ with
$4\le d\le8$, and put $k=d-2$. If \eqref{eq:local-admissibility} holds,
then
\[
 S_{f,k}(X)=c_{f,k}X+o_f(X),\qquad
 c_{f,k}=\prod_p\left(1-\frac{\rho_f(p^k)}{p^k}\right)>0.
\]
No primitivity or sign condition on the leading coefficient is required.
\end{theorem}

For quartics this is a squarefree-value density theorem. In degrees
$5,6,7,8$, the respective excluded powers are $3,4,5,6$.
The statement concerns positive integer inputs and imposes no bound on
the height of the fixed polynomial's coefficients. The error term is not
asserted to be uniform as $f$ varies.

In particular, the quartics $x^4+2$ and $x^4+1$ have positive squarefree-value
density. The first is Eisenstein at $2$; the second is the eighth cyclotomic
polynomial. For $x^4+2$, its value at zero excludes a fixed odd prime square,
and its value at one excludes a fixed square of $2$. The value at zero
of $x^4+1$ is one, so its local condition is immediate.

% [introduction lines 43-109 omitted: history and related work]
The new argument in this paper concerns degrees four through eight.
For higher degrees, Browning's theorem supplies a direct consequence.

\begin{corollary}\label{cor:all-degrees}
Let $f\in\Z[x]$ be irreducible over $\Q$ of degree $d\ge4$, and
put $k=d-2$. If $f$ has no fixed prime $k$th-power divisor, then
\[
 S_{f,k}(X)=c_{f,k}X+o_f(X),\qquad
 c_{f,k}=\prod_p\left(1-\frac{\rho_f(p^k)}{p^k}\right)>0.
\]
\end{corollary}
